\chapter{The Riemann Hypothesis Reformulation}
\label{v4-arith:rh-reformulation}

The zero-placement clause of F-RH is the Riemann Hypothesis.  The full
F-RH assertion also contains the structural assertion that
\(\mathcal{V}^{\mathrm{prim}}\) carries a positive-definite Hermitian
structure compatible with arithmetic Koszul self-duality and arithmetic
Tate's functional equation and explicit formula. The structural theorems of
\Cref{v4-arith:closure} place P1, P2, P3 on their stated domains; the
zero-placement condition is a determinant-divisor and positive-trace
assertion, equivalently a positive Hilbert--P\'olya Li trace after the
regularised trace identity has been supplied.

\section{The Conjecture F-RH}
\label{v4-arith:rh:F-RH-statement}

\begin{conjecture}[arithmetic unitary zero-placement: F-RH]
\label{v4-arith:conj:F-RH}
\(\mathcal{V}^{\mathrm{prim}}\) admits a unique positive-definite
Hermitian structure compatible with arithmetic Koszul self-duality
\(\kappa^{\mathrm{Ar}}+(\kappa^{!})^{\mathrm{Ar}}=0\) and with
Tate's completed-zeta character and explicit formula. Under this structure, the non-trivial
zeros of the entire character \(\xi_{\mathbb R}(s)\) of
\(\mathcal{V}^{\mathrm{prim}}\) lie on the self-dual locus
\(\mathrm{Re}(s)=1/2\).
\end{conjecture}

\begin{remark}[zero-placement and RH]
\label{v4-arith:rh:rem:equivalence-to-RH}
The Tate character supplied by P2 is the meromorphic completed zeta
function
\[
\Lambda_{\zeta}(s):=\pi^{-s/2}\Gamma(s/2)\zeta(s),
\]
with simple poles at \(s=0,1\). The entire Li--Hadamard object is
\[
\xi_{\mathbb R}(s):=\frac12s(s-1)\Lambda_{\zeta}(s).
\]
Here \(\xi\) means \(\xi_{\mathbb R}\) unless
\(\Lambda_{\zeta}\) is written explicitly. The zeros of
\(\xi_{\mathbb R}\) are exactly the non-trivial zeros of \(\zeta\);
the gamma poles cancel the negative even trivial zeros of
\(\zeta\), and the factor \(s(s-1)\) removes the Tate poles.
Thus the zero-placement assertion in F-RH is exactly RH.  The
existence and compatibility of the positive Hermitian structure are
additional arithmetic structure, not a consequence of RH alone.
\end{remark}

\section{Three-Pillar Implication Structure}
\label{v4-arith:rh:three-pillar}

\begin{theorem}[zero-placement from three pillars plus Hilbert--P\'olya--Li]
\label{v4-arith:rh:thm:three-pillar-implication}
\ClaimStatusProvedHere. Suppose:
\begin{itemize}
\item[(P1)] \(\kappa^{\mathrm{Ar}}_{\mathsf{G}}+(\kappa^{!})^{\mathrm{Ar}}_{\mathsf{G}}=0\)
for \(\mathcal{V}^{\mathrm{prim}}_{\mathrm{Heis}}\)
(\Cref{v4-arith:thm:P1-resolution}).
\item[(P2)] \(\xi_{\mathbb R}(s)\) satisfies Tate's functional equation
and the Weil explicit formula, with rank residue \(r_{\zeta}=2\)
(\Cref{v4-arith:thm:P2-resolution}).
\item[(P3)] \(\mathcal{V}^{\mathrm{prim}}\) carries a positive-definite
Petersson inner product compatible with arithmetic Koszul self-duality
(\Cref{v4-arith:thm:P3-resolution}).
\item[(H-P\(^{\mathrm{Ar}}\))] The operator \(\Theta_\xi\) of
\Cref{v4-arith:rh:hyp:HPAr} has determinant divisor
\(\mathrm{div}\,\xi_{\mathbb R}\) and the Weil--Connes trace distribution.
\item[(HP--Li)] The Li transforms of \(\Theta_\xi\) have non-negative
regularised trace.
\end{itemize}
Then the zero-placement clause of F-RH holds. Equivalently, RH holds.
\end{theorem}

\begin{proof}
P1 gives the functional equation \(\xi_{\mathbb R}(s)=\xi_{\mathbb R}(1-s)\).
P2 gives the explicit formula identifying the zero-side distribution
with the prime and archimedean terms. P3 supplies the positive
Hermitian domain. H-P\(^{\mathrm{Ar}}\) places the zero divisor on the
normal operator \(\Theta_\xi\), and HP--Li makes every Li trace
non-negative. Li's theorem then gives RH, which is exactly the
zero-placement clause.
\end{proof}

\begin{remark}[HP--Li as residual condition]
\label{v4-arith:rh:rem:VB-single-open}
With P1, P2, P3 resolved to Theorem or Theorem-conditional,
the zero-placement clause of F-RH reduces to the construction of the
Hilbert--P\'olya zero-divisor operator and to HP--Li positivity. This is
the arithmetic form of the classical Weil--Li positivity problem, not a
formal consequence of Virasoro representation theory.
\end{remark}

\section{Li-Positivity as Arithmetic Virasoro Unitarity}
\label{v4-arith:rh:li-positivity}

\subsection{Li's Criterion}

\begin{theorem}[Li 1997, J.~Number Theory 65]
\label{v4-arith:rh:thm:Li-criterion}
\ClaimStatusProvedElsewhere. RH is equivalent to \(\lambda_{n}\geq 0\)
for all \(n\geq 1\), where
\begin{equation}
\lambda_{n}\;:=\;\sum_{\rho}\biggl(1-\Bigl(1-\frac{1}{\rho}\Bigr)^{n}\biggr)\;=\;\frac{1}{(n-1)!}\frac{d^{n}}{ds^{n}}\bigl[s^{n-1}\log\xi(s)\bigr]\bigg|_{s=1}
\end{equation}
(sum over non-trivial zeros of \(\zeta\), counted with multiplicity).
\end{theorem}

Li's \(\lambda_{n}\) admits a direct Hilbert--P\'olya transform
interpretation. The transform lives on the zero-divisor operator
\(\Theta_{\xi}\), not on the Virasoro/Fock energy operator.

\subsection{Hilbert--P\'olya Transform Computation}

The Li criterion is a scalar positivity statement. To embed it in the
arithmetic chiral-algebra formalism and exploit the Petersson inner
product of P3, one requires an operator \(\Theta_{\xi}\) on the
arithmetic Heisenberg Hilbert completion whose regularised trace
recovers the Li coefficients. The following hypothesis names this
operator precisely.

\begin{hypothesis}[arithmetic Hilbert--P\'olya, H-P\(^{\mathrm{Ar}}\)]
\label{v4-arith:rh:hyp:HPAr}
There exists a Hilbert completion \(\mathcal H_{\xi}\) of the primary
arithmetic Heisenberg core, a dense invariant core
\(\mathcal D_{\xi}\subset\mathcal H_{\xi}\), and a self-adjoint
unbounded operator
\[
H_{\mathrm{HP}}:\mathcal D_{\xi}\subset\mathcal H_{\xi}\longrightarrow
\mathcal H_{\xi}
\]
in the Petersson/GNS inner product. Put
\begin{equation}
\label{v4-arith:rh:def:ThetaXi}
\Theta_{\xi}\;:=\;\frac{1}{2}\,\mathrm{id}\;+\;iH_{\mathrm{HP}}.
\end{equation}
The regularised determinant and trace distribution of
\(\Theta_{\xi}\) satisfy
\begin{equation}
\label{v4-arith:rh:eq:HP-determinant}
\det_{\infty}(s\mathrm{id}-\Theta_{\xi})=\xi_{\mathbb R}(s),
\end{equation}
and, for every admissible Weil--Connes Paley--Wiener test function
\(h\),
\begin{equation}
\label{v4-arith:rh:eq:HP-trace}
\mathrm{Tr}_{\mathrm{reg}}\bigl(h(H_{\mathrm{HP}})\bigr)
=\sum_{\rho}h(\gamma_{\rho}),
\qquad \rho=\frac12+i\gamma_{\rho}.
\end{equation}
Here the notation \(\rho=\tfrac12+i\gamma_{\rho}\) is part of the
hypothesis: self-adjointness of \(H_{\mathrm{HP}}\) together with the
determinant identity forces the zero divisor onto the critical line.
The symbolic kernel \(\sum_{\rho}e^{-t\rho}\) is only the
distributional wave-kernel shorthand obtained from
\eqref{v4-arith:rh:eq:HP-trace} by admissible test-function
regularisation; it is not an ordinary heat trace of a non-negative
self-adjoint operator.
\end{hypothesis}

\begin{remark}[form of H-P\(^{\mathrm{Ar}}\)]
\label{v4-arith:rh:rem:HPAr-form}
H-P\(^{\mathrm{Ar}}\) is the Hilbert--P\'olya conjecture stated
inside the arithmetic chiral-algebra formalism. It is not
derivable from P1, P2, P3 alone: P2 supplies the completed-zeta
functional equation and explicit formula, and P3 supplies the
Petersson domain and inner product on the restricted sph-finite
cuspidal core. The
functional-analytic upgrade to a genuine self-adjoint ordinate
operator \(H_{\mathrm{HP}}\), together with the determinant identity
\eqref{v4-arith:rh:eq:HP-determinant} and the distributional trace
identity \eqref{v4-arith:rh:eq:HP-trace}, is an independent input and
is treated here as an explicit hypothesis.
\end{remark}

\begin{remark}[primary-sector Witt--Virasoro reconciliation]
\label{v4-arith:rh:rem:Virasoro-reconciliation}
\Cref{v4-arith:polyakov} proves that the
\emph{full} \(\mathcal{V}^{\mathrm{prim}}\) carries only a Witt
algebra, not an honest Virasoro: the infinite-tensor central term
\(c_{\mathrm{full}}=\sum_{p}c^{(p)}\) diverges, and no
normal-ordering procedure preserves the Virasoro bracket on the
full tensor. The determinant scalar \(r_{\zeta}=2\)
(\Cref{v4-arith:thm:regularised-c}) exists only as the residue
\(2\cdot\mathrm{Res}_{s=1}\zeta(s)\), not as a trace and not as a
global Virasoro central charge.

The operator \(H_{\mathrm{HP}}\) of H-P\(^{\mathrm{Ar}}\) is therefore
\emph{not} a Virasoro zero-mode on the full
\(\mathcal{V}^{\mathrm{prim}}\). It is a separate spectral
realisation operator. The restriction of the per-prime Virasoro
conformal modes to a primary finite-support core is still useful for
Fock computations, but the claim that these modes extend to a closed
Virasoro representation is itself a substantive conjecture:
\end{remark}

\begin{conjecture}[primary-Virasoro reconciliation]
\label{v4-arith:rh:conj:primary-Virasoro}
On the primary subspace
\(\mathcal{V}^{\mathrm{prim}}_{\mathrm{Heis},\mathrm{prim}}\), the
finite-support conformal modes of \Cref{v4-arith:polyakov} extend
to closed operators \(L_{m}^{\mathrm{PV}}\) satisfying an honest
Virasoro bracket
\begin{equation}
[L_{m}^{\mathrm{PV}},L_{n}^{\mathrm{PV}}]\;=\;(m-n)L_{m+n}^{\mathrm{PV}}+\frac{c_{\mathrm{PV}}}{12}(m^{3}-m)\delta_{m+n,0},\qquad c_{\mathrm{PV}}=2.
\end{equation}
The equality \(c_{\mathrm{PV}}=r_{\zeta}\) is an extra compatibility
condition with the rank residue of
\Cref{v4-arith:thm:regularised-c}. The residue calculation alone does
not construct \(L_{m}^{\mathrm{PV}}\).
\end{conjecture}

\begin{remark}[operator dictionary]
\label{v4-arith:rh:rem:spectral-convention}
Four operators must remain distinct. \(L_{0}^{\mathrm{Vir}}\) is the
Fock/Virasoro conformal-energy operator. \(N_{\mathrm{BC}}\) is the
Bost--Connes Hamiltonian, with eigenvalues \(\log m\) on the
zero-mode lattice and partition function \(\zeta(\beta)\).
\(H_{\mathrm{HP}}\) is the hypothetical self-adjoint determinant-trace
operator. \(\Theta_{\xi}=1/2+iH_{\mathrm{HP}}\) is the normal
zero-divisor operator whose determinant divisor is
\(\mathrm{div}\,\xi_{\mathbb R}\). Conflating
these four objects turns Euler-product positivity into a false
Hilbert--P\'olya argument.
\end{remark}

\begin{remark}[Tate-pole and trivial-zero normalisation]
\label{v4-arith:rh:rem:trivial-zero}
\(\Lambda_{\zeta}(s)=\pi^{-s/2}\Gamma(s/2)\zeta(s)\) is meromorphic,
not entire: the negative even trivial zeros of \(\zeta\) are cancelled
by gamma poles, and \(\Lambda_{\zeta}\) has poles at \(0\) and \(1\).
The entire function
\(\xi_{\mathbb R}(s)=\frac12s(s-1)\Lambda_{\zeta}(s)\) removes these
Tate poles and has as zeros exactly the non-trivial zeros. Thus
H-P\(^{\mathrm{Ar}}\) is a statement about the primary determinant
\(\xi_{\mathbb R}\), while the meromorphic Tate character
\(\Lambda_{\zeta}\) belongs to the cohomological/superdeterminant
surface.
\end{remark}

\begin{proposition}[Li coefficient = Hilbert--P\'olya transform trace]
\label{v4-arith:rh:prop:Li-is-Virasoro}
\ClaimStatusProvedHere \emph{(conditional on
\Cref{v4-arith:rh:hyp:HPAr})}.
Assume H-P\(^{\mathrm{Ar}}\). For \(n\geq1\), define the normal
Li-transform operator
\begin{equation}
\label{v4-arith:rh:eq:Li-transform-operator}
\mathcal L_{n}:=\mathrm{id}-\bigl(\mathrm{id}-\Theta_{\xi}^{-1}\bigr)^{n}.
\end{equation}
Then
\begin{equation}
\label{v4-arith:rh:eq:Li-transform-trace}
\lambda_{n}=\mathrm{Tr}_{\mathrm{reg}}(\mathcal L_{n}).
\end{equation}
This is a trace identity for the Hilbert--P\'olya zero-divisor operator
\(\Theta_{\xi}\), not a polynomial identity in
\(L_{0}^{\mathrm{Vir}}\) or \(N_{\mathrm{BC}}\).
\end{proposition}

\begin{proof}
The determinant identity
\(\det_{\infty}(s\mathrm{id}-\Theta_{\xi})=\xi_{\mathbb R}(s)\)
identifies the determinant divisor of \(\Theta_{\xi}\), in the
regularised determinant sense, with the non-trivial zero divisor.
Applying the holomorphic functional calculus to
\(f_{n}(z)=1-(1-z^{-1})^{n}\) gives
\[
\mathrm{Tr}_{\mathrm{reg}}(f_{n}(\Theta_{\xi}))
=\sum_{\rho}\left(1-\left(1-\frac1{\rho}\right)^{n}\right),
\]
which is Li's coefficient by \Cref{v4-arith:rh:thm:Li-criterion}.
\end{proof}

\begin{corollary}[Li positivity from positive HP trace]
\label{v4-arith:rh:cor:VB-implies-Li}
\ClaimStatusProvedHere \emph{(conditional on P2, P3, H-P\(^{\mathrm{Ar}}\))}.
Assume the regularised Hilbert--P\'olya trace is positive on the
Weil positive cone and that the Li transforms
\(\mathcal L_{n}\) of \eqref{v4-arith:rh:eq:Li-transform-operator}
belong to the closure of that positive cone after the standard
Li/Weil approximation. Then \(\lambda_{n}\geq0\) for every
\(n\geq1\). The converse is \emph{not} an operator-positivity theorem.
\end{corollary}

\begin{proof}
By the positivity assumption,
\(\mathrm{Tr}_{\mathrm{reg}}(\mathcal L_{n})\geq0\). The trace
identity \eqref{v4-arith:rh:eq:Li-transform-trace} identifies this
number with \(\lambda_{n}\). The reverse implication would require a
determinate positive trace representation of the full zero measure
from the scalar Li sequence; that is the classical Li/Weil positivity
problem, not a consequence of P1--P3.
\end{proof}

\begin{theorem}[positive Hilbert--P\'olya trace implies zero-placement]
\label{v4-arith:rh:thm:VB-implies-F-RH}
\ClaimStatusProvedHere \emph{(conditional on P2, P3, H-P\(^{\mathrm{Ar}}\))}.
Assume P2, P3, and the determinant and regularised trace identities
of H-P\(^{\mathrm{Ar}}\). If HP--Li positivity holds:
\begin{itemize}
\item[\textbf{(HP--Li)}] \emph{Positive Hilbert--P\'olya trace.}
\(\mathrm{Tr}_{\mathrm{reg}}(\mathcal L_{n})\geq0\) for every
\(n\geq1\), where \(\mathcal L_{n}\) is
\eqref{v4-arith:rh:eq:Li-transform-operator},
\end{itemize}
then the zero-placement clause of F-RH holds. Equivalently, RH holds.
If the self-adjointness assertion in H-P\(^{\mathrm{Ar}}\) is also
assumed, then the determinant identity already gives zero-placement;
the scalar Li route uses only the determinant and trace identities.
\end{theorem}

\begin{proof}
HP--Li gives \(\lambda_{n}\geq0\) for all \(n\geq1\) by
\Cref{v4-arith:rh:prop:Li-is-Virasoro}. Li positivity is equivalent
to RH (Li 1997,
\Cref{v4-arith:rh:thm:Li-criterion}); this is exactly the
zero-placement clause of F-RH
(\Cref{v4-arith:rh:rem:equivalence-to-RH}).
\end{proof}

\begin{remark}[one-directional implication and operator discipline]
\label{v4-arith:rh:rem:one-directional}
The polynomial \((L_{0}^{\mathrm{Ar}})^{n}-(L_{0}^{\mathrm{Ar}}-\mathrm{id})^{n}\)
is \emph{not} Li's transform; \(L_{0}^{\mathrm{Ar}}\) would need to
serve simultaneously as Fock energy, Bost--Connes Hamiltonian, and
determinant-trace operator, three incompatible roles (see
\Cref{v4-arith:rh:rem:spectral-convention}). The correct statement
is the one-directional HP--Li positivity
\(\Rightarrow\) zero-placement, where H-P\(^{\mathrm{Ar}}\) supplies
\(\Theta_{\xi}\). The converse is the assertion that zero-placement
supplies a positive Weil trace representation of the full zero
distribution; this is the Hilbert--P\'olya/Connes--Deninger
positivity problem.
\end{remark}

\begin{remark}[content of F-RH as a reformulation]
\label{v4-arith:rh:rem:what-is-gained}
F-RH separates RH from the extra arithmetic structures carried by
\(\mathcal{V}^{\mathrm{prim}}\).
\begin{enumerate}
\item The spectral object is made precise: the self-adjoint
determinant form of H-P\(^{\mathrm{Ar}}\) is a Hilbert--P\'olya
realisation of RH, while HP--Li positivity is Li's scalar
RH-equivalent positivity written through the zero-divisor operator after the
trace identity is fixed.
\item A structural constraint on zero-placement emerges: the
Petersson positivity of P3 propagates through the Weil/Connes trace
functional to impose positivity of the Li transforms, provided
the trace representation is determinate.
\item The reformulation connects RH mechanisms (Connes, Deninger,
Meyer) to positive-trace and factorisation methods, with the
arithmetic explicit formula replacing compact-CFT modularity.
\item Admissible approximants to the Li/Weil test functions can be
evaluated against the Bost--Connes and Hecke trace formulae.
Positivity is a
distributional trace question, not finite Fock positivity of
\(T_{n}\).
\end{enumerate}
\end{remark}

\section{Comparison to Existing RH Mechanisms}
\label{v4-arith:rh:comparison}

\begin{table}[h]
\centering
\small
\begin{tabular}{l|l|l}
Mechanism & Input & Analytic theorem \\
\hline
Hilbert--P\'olya (folklore) & Self-adjoint op.~with determinant divisor \(\mathrm{div}\,\xi\) & Operator not exhibited \\
Connes (adele class space) & KMS\({}_{\beta=1}\) Bost--Connes positivity & Trace-formula positivity \\
Deninger (\(H^{\bullet}_{\Delta}\)) & Purity of Frobenius on arith.~\(H^{1}\) & Cohomology not constructed \\
Meyer (Schwartz space) & Mellin self-adjoint operator plus missing zero-divisor quotient & Operator not exhibited \\
Arakelov height-zeta & Height-function positivity & Equivalence to RH not established \\
Weil positivity & Positivity of Weil inner product & Known equivalent, no proof \\
Li positivity & \(\lambda_{n}\geq 0\) for all \(n\geq 1\) & Known equivalent, no proof \\
Riemann--von Mangoldt & Explicit formula approach & Approaches RH asymptotically \\
Montgomery--Dyson & GUE statistics of zeros & Statistical, not deterministic \\
\textbf{Vol~IV zero-placement} & \textbf{H-P\(^{\mathrm{Ar}}\) + positive trace} & \textbf{Self-adjoint zero-divisor operator not constructed} \\
\end{tabular}
\caption{The zero-placement clause placed against existing RH mechanisms.}
\label{v4-arith:rh:table:comparison}
\end{table}

\begin{remark}[relations]
\label{v4-arith:rh:rem:comparison-relations}
\begin{enumerate}
\item HP--Li positivity is Li positivity written on the
Hilbert--P\'olya zero-divisor operator \(\Theta_{\xi}\). It implies
zero-placement, hence RH, via
\Cref{v4-arith:rh:thm:VB-implies-F-RH}. A Virasoro polynomial
\(T_{n}\) is not a substitute for this transform.
\item The zero-placement clause is governed by Weil positivity:
Weil's inner product on test functions corresponds to the Petersson
inner product on \(\mathcal{V}^{\mathrm{prim}}\) via Tate's restricted
tensor product. F-RH adds the arithmetic Hermitian, Koszul, and
completed-zeta trace structures around this clause.
\item F-RH embeds Hilbert--P\'olya explicitly as H-P\(^{\mathrm{Ar}}\)
(\Cref{v4-arith:rh:hyp:HPAr}): it asks for a self-adjoint operator
\(A\), the positive form from P3, the determinant divisor of
\(\Theta_{\xi}=1/2+iA\), the
regularised non-trivial-zero trace distribution, and
\(\mathbb{Z}/2\)-symmetry under the arithmetic
\(S\)-modular involution from P1. F-RH does not replace
Hilbert--P\'olya: H-P\(^{\mathrm{Ar}}\) is a separate spectral
hypothesis in the sufficient implication below.
\item F-RH is closer to Connes than to Deninger: both Connes and
F-RH use unitarity plus Bost--Connes positivity; Deninger
postulates a cohomology theory that F-RH does not need. The
difference from Connes is that F-RH carries the positivity
through arithmetic holomorphic-topological (HT) observables and a
Weil-positive trace rather than through the bare trace-formula construction.
\end{enumerate}
\end{remark}

\section{Residual Conditions}
\label{v4-arith:rh:residual}

After the structural reduction of \Cref{v4-arith:closure} and the
F-RH reformulation above, the RH-relevant part is controlled by the
following conditions.  The spectral conditions imply zero-placement;
the categorical conditions construct the trace object on which that
problem is stated.

\begin{enumerate}
\item[\textbf{(HP--Li)}] \emph{Positive Hilbert--P\'olya trace.}
\(\mathrm{Tr}_{\mathrm{reg}}(\mathcal L_{n})\geq0\) for every
\(n\geq1\), where \(\mathcal L_{n}\) is the Li transform of
\(\Theta_{\xi}\)
(\Cref{v4-arith:rh:thm:VB-implies-F-RH}).
\emph{Sufficient for RH given P2, P3, H-P\(^{\mathrm{Ar}}\).}
\item[\textbf{(H-P\(^{\mathrm{Ar}}\))}] \emph{Arithmetic Hilbert--P\'olya.}
There is a self-adjoint operator \(A\) such
that \(\Theta_{\xi}=1/2+iA\) has determinant divisor
\(\mathrm{div}\,\xi_{\mathbb R}\) and Weil--Connes trace distribution equal to the
non-trivial-zero distribution
(\Cref{v4-arith:rh:hyp:HPAr}).
\emph{Required for the HP--Li \(\Rightarrow\) zero-placement chain; not
derivable from P1, P2, P3.}
\item[(arith.~Positselski)] Arithmetic Positselski equivalence
(\Cref{v4-arith:thm:arithmetic-positselski}).
\emph{Required for full P1 beyond the forward implication.}
\item[(Koszul--Pet compat)] Identification of arithmetic Koszul
pairing with adelic Petersson pairing
(\Cref{v4-arith:cl:cj:koszul-Petersson-compat}).
\emph{Required for P3.}
\item[(\(\mathrm{L}_{\mathrm{ACD}}\))] Adelic categorical descent
for \(\mathrm{D}^{\mathrm{sph\text{-}fin}}\)
(\Cref{v4-arith:cl:cor:F1-single-lemma}).
\emph{Required for F1, hence for six-window placement.}
\item[(F3.BZN\(^{\mathrm{Ar}}\))] Arithmetic Ben-Zvi--Nadler
identification (\Cref{v4-arith:thm:F3-reduction}).
\emph{Required for F3, hence for the categorical trace presentation
of \(\mathcal{V}^{\mathrm{prim}}\).}
\end{enumerate}

Items (arith.~Positselski), (Koszul--Pet compat), H-P\(^{\mathrm{Ar}}\),
and HP--Li positivity lie on the P1--P2--P3 branch of the
zero-placement clause and together form a sufficient set for RH in the
chiral-algebra formalism. Items (\(\mathrm{L}_{\mathrm{ACD}}\)) and
(F3.BZN\(^{\mathrm{Ar}}\)) are independent of RH proper and lie on
the six-window presentation; they would state
\(\mathcal{V}^{\mathrm{prim}}\) as a theorem-grade object but do
not imply RH.

\section{A Sufficient Path to RH}
\label{v4-arith:rh:sufficient-path}

\begin{theorem}[sufficient implication]
\label{v4-arith:rh:thm:sufficient-path}
\ClaimStatusProvedHere. RH follows from the following implication in
the chiral-algebra formalism:
\begin{center}
\begin{tikzcd}[row sep=1em, column sep=small]
\text{P1} \ar[r, "\text{arith.~Positselski}"] & \text{P1 functional equation}\\
\text{P2} \ar[r, "\text{Tate 1950}"] & \text{P2 unconditional}\\
\text{P3} \ar[r, "\text{Koszul--Pet compat}"] & \text{P3 restricted unitarity}\\
\text{(P1, P2, P3, H-P}^{\mathrm{Ar}}\text{)} \ar[r, "\text{HP--Li}"] & \text{zero-placement} \ar[r, "\Leftrightarrow"] & \text{RH}
\end{tikzcd}
\end{center}
Four inputs suffice: arithmetic Positselski, Koszul--Pet
compatibility, arithmetic Hilbert--P\'olya (H-P\(^{\mathrm{Ar}}\)),
and positive Hilbert--P\'olya Li trace (HP--Li).
\end{theorem}

\begin{proof}
P2 is the completed-zeta character/explicit-formula input
(\Cref{v4-arith:thm:P2-resolution}); P1 is the functional-equation
branch conditional on arithmetic Positselski
(\Cref{v4-arith:thm:P1-resolution}); P3 is the positive-form input
conditional on the restricted Koszul--Petersson compatibility
(\Cref{v4-arith:thm:P3-resolution}). Given these inputs plus
H-P\(^{\mathrm{Ar}}\),
\Cref{v4-arith:rh:thm:VB-implies-F-RH} gives
HP--Li \(\Rightarrow\) zero-placement, and zero-placement
\(\Leftrightarrow\) RH via Li's
criterion (\Cref{v4-arith:rh:rem:equivalence-to-RH},
\Cref{v4-arith:rh:thm:Li-criterion}).
\end{proof}

\begin{remark}[analytic content of the implication]
\label{v4-arith:rh:rem:sufficient-path-paper}
\begin{description}
\item[arith.~Positselski] non-formal theorem; Iwasawa-theoretic Positselski.
\item[Koszul--Pet compat] non-formal theorem; compatibility of the
arithmetic Koszul and adelic Petersson pairings.
\item[H-P\(^{\mathrm{Ar}}\)] This is the Hilbert--P\'olya
operator-construction problem in the arithmetic Fock setting; the
domain theorem alone does not supply the zero-divisor determinant.
\item[HP--Li] This is the positive-trace form of the
classical Hilbert--P\'olya / Weil positivity problem. If
tractable through Li coefficients, the
\(n=2\) case (\(\lambda_{2}\geq 0\)) has been
computed numerically to \(n\sim 10^{5}\) zeros of \(\zeta\) (Coffey
2005; Voros 2003 for the first few). The assertion for all \(n\) is
Li's RH-equivalent positivity theorem.
\end{description}
The arithmetic formalism does not shorten the Hilbert--P\'olya
obstacle, but it places the positivity in a precise chain-level trace
setting where factorisation techniques may be applied without
conflating them with the zero-divisor theorem.
\end{remark}

\section{Relation to the Arithmetic Branch Essential Theorem}
\label{v4-arith:rh:load-bearing-relation}

The arithmetic branch's essential theorem
(\Cref{v4-arith:thm:G-row-Riemann-functional-equation},
\Cref{v4-arith:thm:P1-resolution}) is
\begin{equation}
\kappa^{\mathrm{Ar}}_{\mathsf{G}}+(\kappa^{!})^{\mathrm{Ar}}_{\mathsf{G}}=0\;\Longleftrightarrow\;\xi(s)=\xi(1-s).
\end{equation}
This gives the functional equation (geometric self-duality of
\(\xi_{\mathbb R}\)), not RH (spectral zero-placement). RH is
strictly beyond the essential theorem, requiring the additional
positive Hilbert--P\'olya trace input.

F-RH asks whether Koszul self-duality, which gives the functional
equation, combined with the positivity of P3 and the explicit-formula
content of P2, suffices to deliver the spectral zero-placement that RH
asserts.
The implication in the chiral-algebra formalism is conditional
on H-P\(^{\mathrm{Ar}}\) and HP--Li positivity. Those two statements
are the Hilbert--P\'olya form of RH native to this formalism.

\section{The F-RH Reduction}
\label{v4-arith:rh:summary}

\begin{enumerate}
\item The zero-placement clause of F-RH
(\Cref{v4-arith:conj:F-RH}) is exactly RH
(\Cref{v4-arith:rh:rem:equivalence-to-RH}). F-RH
also asks for the positive Hermitian structure compatible with
arithmetic Koszul self-duality, Tate's functional equation, and the
explicit formula.
\item The zero-placement clause is implied by H-P\(^{\mathrm{Ar}}\) and
HP--Li positivity together with P1 + P2 + P3
(\Cref{v4-arith:rh:thm:three-pillar-implication}).
\item The zero-placement clause is implied by HP--Li positivity
(positive Hilbert--P\'olya trace;
\Cref{v4-arith:rh:thm:VB-implies-F-RH}) via Li's criterion, given
the determinant and trace identities of H-P\(^{\mathrm{Ar}}\). If the
full determinant-trace H-P\(^{\mathrm{Ar}}\) hypothesis is assumed,
zero-placement is already part of the determinant-trace realisation.
The converse direction from zero-placement to HP--Li is the
positive-trace realisation problem.
\item Li's coefficients \(\lambda_{n}\) are Hilbert--P\'olya
transform traces of \(\Theta_{\xi}\)
(\Cref{v4-arith:rh:prop:Li-is-Virasoro}).
\item Li positivity and Weil positivity govern the zero-placement
clause. F-RH joins this clause to the arithmetic Hermitian,
Koszul, and completed-zeta trace structures.
\item The sufficient implication to RH uses four inputs:
arith.~Positselski, Koszul--Pet compat, H-P\(^{\mathrm{Ar}}\), and
HP--Li positivity
(\Cref{v4-arith:rh:thm:sufficient-path}). This is the
\emph{four-input conditional conjunction} stated in
\Cref{v4-arith:w3-comparison-B} and made precise as
\Cref{v4-arith:w3-comparison-B:thm:four-input-branch}: each input
carries its own analytic content (five Positselski sub-items,
KMS\(_{1}\)/GNS identification beyond the sph-fin cuspidal core,
zeta-zero determinant-trace matching, RH-equivalent positive trace),
and the conjunction implies zero-placement.
\item RH is strictly beyond the branch's essential theorem
(functional equation), and HP--Li positivity is the
Li positive-trace form of RH within the chiral-algebra formalism.
The operator dictionary of
\Cref{v4-arith:rh:rem:one-directional} preserves only the
forward implication HP--Li\(\Rightarrow\)zero-placement; the converse is the
positive determinant-trace realisation problem
(\Cref{v4-arith:hpvbconv:cor:three-strictly-finer}) and does not
help the forward direction.
\end{enumerate}
